% Part: applied-modal-logics
% Chapter: epistemic-logic
% Section: introduction

\documentclass[../../../include/open-logic-section]{subfiles}

\begin{document}

\olfileid{aml}{el}{int}

\olsection{سريزه}

لکه څنګه چې مودالي منطق له \emph{مودالي قضيو} او د هغو تر منځ د لازم راتلو
له اړيکو سره کار لري، معرفتي منطق له \emph{معرفتي قضيو} او د هغو تر منځ
د لازم راتلو له اړيکو سره کار لري۔ دلته يو-ځاييز نښلوونکي د امکان او
لزوم د ښودلو پر ځاے د پوهې يا باور له مخې تفسيرېږي، څو پوهه او باور
مدل کړي۔ د بېلګې په توګه، ښايي لاندې دعوې څرګندول وغواړو:

\begin{enumerate}
\item ريچارډ پوهېږي چې کالګري په البرټا کښې ده۔
\item آډري فکر کوي چې شونې ده يو سپے پر صوفه وي۔
\item ريچارډ پوهېږي چې آډري پوهېږي د هغې ټولګے د نهې په ورځو وي۔
\item هر څوک پوهېږي چې يو کال ۱۲ مياشتې لري۔
\end{enumerate}
اوسنے معرفتي منطق ډېر ځله د يااکو هينتيکا (Jaako Hintikka) د ۱۹۶۲ کال
کتاب \emph{پوهه او باور} (Knowledge and Belief) ته منسوبېږي۔ دا کتاب
په داسې وخت کښې ليکل شوے و چې په منطق کښې د ممکنو نړۍګانو معنٰی پوهنه
لا ډېره کارېده۔ په حقيقت کښې معرفتي منطقونه تر ډېره د نورو مودالي
منطقونو هماغه معنايي وسيلې کاروي، خو بل ډول ئې تفسيروي۔ اصلي بدلون
په دې کښې دے چې \emph{د لاسرسي اړيکه} د څۀ نماينده ګڼو۔ په معرفتي
منطقونو کښې دا اړيکه د \emph{معرفتي امکان} يوه بڼه ښيي۔ وبه ګورو چې
مدل کېدونکے معرفتي تصور د لاسرسي پر اړيکه زموږ د مطلوبو شرطونو په
ټاکلو څۀ اغېز کوي۔ همدارنګه به وګورو چې د مطابقت تيوري ته د معرفتي
تفسير په ورکولو سره څۀ پېښېږي۔ په پورته بېلګو کښې دوه عاملان، ريچارډ
او آډري، ياد شوي او د هغو څيزونو تر منځ اړيکه هم شته چې هر يو ئې
پېژني۔ هغه معرفتي منطقونه چې دلته ئې څېړو څو-عامله منطقونه به وي؛
په هغو کښې دا ډول خبرې څرګندېدے شي۔ برعکس، يو-عامله معرفتي منطق
يوازې د يوه کس د پوهې يا باور په اړه خبرې کوي۔

\end{document}
